\documentclass[12pt, a4paper]{article}
\usepackage[utf8]{inputenc}
\usepackage[T1]{fontenc}
\usepackage[french]{babel}
\usepackage{amsmath, amssymb, amsthm}
\usepackage{geometry}
\usepackage{listings}
\geometry{margin=2.5cm}

\newtheorem{theorem}{Th\'eor\`eme}[section]
\newtheorem{lemma}[theorem]{Lemme}
\newtheorem{definition}[theorem]{D\'efinition}
\newtheorem{remark}[theorem]{Remarque}

\title{Sur la Conjecture d'Erd\H{o}s-Sierpi\'nski : Analyse Alg\'ebrique et D\'ecomposition Modulaire}
\author{Charles EDOU NZE}
\date{\today}

\lstdefinelanguage{lean}{
  keywords={import, def, theorem, lemma, by, admit, Prop, Nat, open, section, Exists, fun, Set, Finset, Real, Filter, Topology},
  sensitive=true,
  comment=[l]--
}

\begin{document}
\maketitle
\thispagestyle{empty}


\begin{abstract}
Ce document traite formellement de la conjecture d'Erd\H{o}s-Sierpi\'nski, qui affirme que pour tout entier $n \ge 2$, l'\'equation $\frac{5}{n} = \frac{1}{x} + \frac{1}{y} + \frac{1}{z}$ poss\`ede des solutions enti\`eres positives $x, y, z$. La structure englobe des d\'efinitions axiomatiques strictes, une revue de la litt\'erature pertinente ax\'ee sur les cribles modulaires et la param\'etrisation des diviseurs, une d\'ecomposition en classes de congruence modulo $5$ et $4$, ainsi que des d\'erivations d\'etaill\'ees couvrant ces classes. Une architecture formelle en Lean 4 est ensuite esquiss\'ee.
\vfill
\noindent \textit{Charles EDOU NZE, chercheur ind\'ependant}
\end{abstract}

\tableofcontents
\newpage

\section{D\'efinitions Axiomatiques}

\begin{definition}
Soit $n \in \mathbb{N}$ avec $n \ge 2$. Un triplet d'entiers positifs $(x, y, z) \in (\mathbb{N}_{>0})^3$ est d\'efini comme une solution \`a l'\'equation Diophantienne d'Erd\H{o}s-Sierpi\'nski s'il satisfait :
\[ \frac{5}{n} = \frac{1}{x} + \frac{1}{y} + \frac{1}{z} \]
\end{definition}

\begin{definition}
Nous d\'efinissons le pr\'edicat $\mathcal{E}(n)$ comme \'etant vrai si et seulement s'il existe un triplet solution $(x, y, z) \in (\mathbb{N}_{>0})^3$ pour un entier donn\'e $n \ge 2$.
\end{definition}

\section{Litt\'erature Contextuelle}

La conjecture d'Erd\H{o}s-Sierpi\'nski formule une question fondamentale concernant la repr\'esentation des nombres rationnels comme sommes de fractions unitaires, g\'en\'eralisant l'\'equation bien connue d'Erd\H{o}s-Straus $\frac{4}{n} = \frac{1}{x} + \frac{1}{y} + \frac{1}{z}$. Des avanc\'ees significatives se sont appuy\'ees sur l'analyse de l'\'equation via l'arithm\'etique modulaire. Les approches impliquant des cribles modulaires ont v\'erifi\'e empiriquement la conjecture jusqu'\`a des bornes massives. Les \'etudes th\'eoriques emploient intensivement des param\'etrisations de diviseurs et des identit\'es polynomiales. Un r\'esultat analytique de premier plan \'etablit que le nombre de solutions est born\'e de mani\`ere poly-logarithmique en moyenne. La m\'ethodologie d\'evelopp\'ee ici exploite ces identit\'es polynomiales param\'etr\'ees par des classes de congruence.

\section{Isolation des Lemmes et Preuves Z\'ero Ellipse}

La r\'eduction de la conjecture aux valeurs premi\`eres $n = p$ est standard, car une solution pour un facteur premier $p$ peut \^etre \'etendue aux nombres compos\'es via la multiplication scalaire. Nous \'etablissons des lemmes sp\'ecifiques pour des classes de congruence distinctes.

\begin{lemma}
Pour tout entier $k \ge 0$, si $n = 5k + 4$, alors $\mathcal{E}(n)$ est vrai.
\end{lemma}

\begin{proof}
Soit $n = 5k + 4$ pour $k \in \mathbb{N}$. Nous consid\'erons le d\'eveloppement en fraction rationnelle de $\frac{5}{5k+4}$. Nous d\'efinissons le triplet de variables $(x, y, z)$ comme suit :
\[ x = k + 1 \]
\[ y = 2(5k+4)(k+1) \]
\[ z = 2(5k+4)(k+1) \]
Nous \'evaluons la somme des fractions unitaires r\'eciproques correspondant \`a cette affectation :
\[ \frac{1}{x} + \frac{1}{y} + \frac{1}{z} = \frac{1}{k+1} + \frac{1}{2(5k+4)(k+1)} + \frac{1}{2(5k+4)(k+1)} \]
Puisque la deuxi\`eme et la troisi\`eme fraction sont identiques, nous les sommons :
\[ \frac{1}{2(5k+4)(k+1)} + \frac{1}{2(5k+4)(k+1)} = \frac{2}{2(5k+4)(k+1)} = \frac{1}{(5k+4)(k+1)} \]
En substituant cela dans la somme originale :
\[ \frac{1}{k+1} + \frac{1}{(5k+4)(k+1)} \]
Pour combiner ces fractions, nous identifions un d\'enominateur commun, qui est $(5k+4)(k+1)$. Nous multiplions le num\'erateur et le d\'enominateur de la premi\`ere fraction par $(5k+4)$ :
\[ \frac{5k+4}{(5k+4)(k+1)} + \frac{1}{(5k+4)(k+1)} = \frac{5k+5}{(5k+4)(k+1)} \]
Nous factorisons le num\'erateur par $5(k+1)$ :
\[ \frac{5(k+1)}{(5k+4)(k+1)} \]
Puisque $k \ge 0$, le terme $(k+1)$ est non nul, ce qui nous permet de l'annuler du num\'erateur et du d\'enominateur :
\[ \frac{5}{5k+4} = \frac{5}{n} \]
Cela d\'emontre que $x = k+1$, $y = 2(5k+4)(k+1)$, et $z = 2(5k+4)(k+1)$ constituent un triplet d'entiers solution valide pour $n = 5k+4$. Puisque $k \ge 0$, tous les d\'enominateurs sont des entiers strictement positifs.
\end{proof}

\begin{lemma}
Pour tout entier $k \ge 1$, si $n = 4k + 3$, alors $\mathcal{E}(n)$ est vrai.
\end{lemma}

\begin{proof}
Soit $n = 4k + 3$ pour $k \in \mathbb{N}_{>0}$.
Nous d\'efinissons $x = k+1$.
La diff\'erence est \'evalu\'ee comme :
\[ \frac{5}{4k+3} - \frac{1}{k+1} = \frac{5(k+1) - (4k+3)}{(4k+3)(k+1)} = \frac{5k+5-4k-3}{(4k+3)(k+1)} = \frac{k+2}{(4k+3)(k+1)} \]
Nous exigeons que $\frac{k+2}{(4k+3)(k+1)}$ soit repr\'esent\'e par $\frac{1}{y} + \frac{1}{z}$.
Nous utilisons l'identit\'e alg\'ebrique $\frac{A+B}{C} = \frac{A}{C} + \frac{B}{C}$. Soit $A = k+1$ et $B = 1$. Alors $A+B = k+2$.
\[ \frac{k+2}{(4k+3)(k+1)} = \frac{k+1}{(4k+3)(k+1)} + \frac{1}{(4k+3)(k+1)} \]
En simplifiant la premi\`ere fraction par annulation du facteur commun $(k+1)$ :
\[ \frac{k+1}{(4k+3)(k+1)} = \frac{1}{4k+3} \]
Ainsi, la diff\'erence se d\'ecompose exactement en deux fractions unitaires :
\[ \frac{1}{4k+3} + \frac{1}{(4k+3)(k+1)} \]
Nous assignons $y = 4k+3$ and $z = (4k+3)(k+1)$. Puisque $k \ge 1$, $x=k+1$, $y=4k+3$, et $z=(4k+3)(k+1)$ sont tous des entiers strictement positifs.
L'identit\'e compl\`ete est :
\[ \frac{5}{4k+3} = \frac{1}{k+1} + \frac{1}{4k+3} + \frac{1}{(4k+3)(k+1)} \]
Ceci conclut la preuve pour cette classe de congruence.
\end{proof}

\section{Architecture pour l'Autoformalisation}
La v\'erification formelle des lemmes susmentionn\'es peut \^etre impl\'ement\'ee dans Lean 4. Les types fondamentaux sont sp\'ecifi\'es comme suit :

Les d\'efinitions et types suivants pr\'esentent une structure pour traduire le probl\`eme et les identit\'es modulaires \'etablies en Lean 4.

\begin{lstlisting}[language=lean, basicstyle=\ttfamily\small, breaklines=true]
import Mathlib.Data.Nat.Basic
import Mathlib.Data.Rat.Basic

open Nat

-- D\'efinition axiomatique de la propri\'et\'e d'Erd\H{o}s-Sierpi\'nski pour un entier n
def SatisfiesErdosSierpinski (n : Nat) : Prop :=
  Exists (fun x => Exists (fun y => Exists (fun z =>
    x > 0 /\ y > 0 /\ z > 0 /\
    (5 : Rat) / n = (1 : Rat) / x + (1 : Rat) / y + (1 : Rat) / z)))

-- Lemme pour la classe de congruence n = 5k + 4
lemma erdos_sierpinski_mod_5_4 (k : Nat) : SatisfiesErdosSierpinski (5 * k + 4) := by
  admit

-- Lemme pour la classe de congruence n = 4k + 3
lemma erdos_sierpinski_mod_4_3 (k : Nat) (hk : k >= 1) : SatisfiesErdosSierpinski (4 * k + 3) := by
  admit

-- \'Enonc\'e g\'en\'eral de la conjecture d'Erd\H{o}s-Sierpi\'nski
theorem erdos_sierpinski_conjecture (n : Nat) (hn : n >= 2) : SatisfiesErdosSierpinski n := by
  admit
\end{lstlisting}

\section*{R\'ef\'erences}
\begin{itemize}
    \item Webb, W. A. (1970). \textit{On $4/n = 1/x + 1/y + 1/z$}. Proceedings of the American Mathematical Society, 25(3), 578-584.
    \item Vaughan, R. C. (1970). \textit{On a problem of Erd\H{o}s, Straus and Schinzel}. Mathematika, 17(2), 193-198.
    \item Sierpi\'nski, W. (1956). \textit{Sur les d\'ecompositions de nombres rationnels en fractions primaires}. Mathesis, 65, 16-32.
\end{itemize}

\end{document}
